% 20-conclusion.tex: the complete mental model, answered.
\section{Conclusion: the whole organism}
\label{sec:conclusion}

A reader who has followed the rings outward should now be able to answer the questions this paper
set out to make answerable. Each answer below is one sentence, and each points back to the section
that carries the mechanism.

\begin{description}
  \item[What the central credential is.] One signed claim, \emph{this authority issued this token to
    this person}, held as a file the holder controls, in a lifecycle with exactly one active state
    (Section~\ref{sec:core}).
  \item[Why the schema surrounds it.] Because a promise enforced only by the people who run the
    system is not a promise; triggers, check constraints and unique indexes bind every client and
    every restore (Section~\ref{sec:walls}).
  \item[How cryptographic trust works.] Post-quantum signatures over digests, checked by two
    independent implementations at issuance and sampled at use, with keys that have a recorded
    lifecycle (Section~\ref{sec:signatures}).
  \item[How verification differs from authorization.] Authenticity is permanent and cacheable;
    authorization is fresh and revocable, decided on the primary or by a short-lived signed assertion
    whose window is the honest cost of going offline (Sections~\ref{sec:verifiers}
    and~\ref{sec:authorization}).
  \item[How privacy is bounded.] By what is stored, not by who may read it: no token identifier on
    a zero-knowledge event, a membership proof that reveals nothing but membership, bounded reads,
    and a stated limitation on cross-verifier correlation (Section~\ref{sec:privacy}).
  \item[How relying parties use credentials.] Through a wallet protocol, a detached verifier, two
    SDKs and a login broker that keeps no record of who logged in where (Section~\ref{sec:use}).
  \item[How authorities federate.] Each with its own key and no central store, by directional
    per-context attestations that never form a chain, and by signed feeds a relying party decides
    on offline (Section~\ref{sec:federation}).
  \item[How transparency constrains authority.] Append-only logs with signed heads, monitors that
    prove history only grew, and witnesses that cosign, gossip, and produce a non-repudiable proof
    when a log tells two stories (Section~\ref{sec:transparency}).
  \item[How institutional exchange works.] Through a signed registry, a gateway that authenticates
    by key and authorizes by the responder's own attestation, receipts that commit to hashes and never
    bodies, a timestamp authority that learns nothing, and signatures that outlive their keys
    (Section~\ref{sec:institutions}).
  \item[How Polaris survives failures.] By drilling them: failover, replicas, partitions, backup
    and restore, chaos and rolling deploys, each with a measured number and a ceiling, on every push
    (Section~\ref{sec:operations}).
  \item[What Atlas and Athena do.] Atlas sees what the system is doing in bounded aggregates;
    Athena understands the structure of authority and the constitution as data; neither may model a
    person (Section~\ref{sec:cognition}).
  \item[Why AI agents change the identity problem.] Because the question becomes what kind of actor
    is acting, under whose authority, within what scope, verifiable without learning anything else
    (Section~\ref{sec:agents}).
  \item[How privacy-preserving proof of personhood might fit.] On the existing uniqueness constraint
    and membership proof, with a scoped nullifier, a holder-side key and a presentation that carries no
    stable handle, none of which is built (Sections~\ref{sec:agents} and~\ref{sec:research}).
  \item[How delegated agent authority might work.] As a signed, scoped, expiring, separately revocable
    grant from a human's credential to an agent, verified by a service that learns only the five
    answers it needs (Section~\ref{sec:agents}).
  \item[How stigmergic debugging could help.] As an outside colony of specialised agents leaving
    evidence on a shared landscape, with confidence drawn only from reproductions, feeding the loop by
    which a reproducible failure becomes a permanent invariant (Section~\ref{sec:research}).
  \item[What Polaris can do today.] Issue, sign, verify, revoke, present, federate, log, witness,
    exchange, timestamp and sign documents, all on notional data, all drilled in continuous
    integration, between instances of one code base (Appendix~\ref{app:status}).
  \item[What it cannot do today.] Hold real identity data, run on hardware tokens or a hardware
    module, prove independence of institutions or implementations, prevent cross-verifier
    correlation, or claim any property no outside party has tested (Section~\ref{sec:limitations}).
  \item[What it is ultimately trying to become.] Identity infrastructure in which no important claim
    has to be taken on the word of a central application, because authority, state, history, privacy
    and the system's own behaviour are exposed to independent verification and constrained by layers
    that do not share each other's assumptions.
\end{description}

The town's register is small. The walls around it are what make it safe to keep. And the reason the
system has so many rings is not that its builders enjoyed rings: each one was added when the ring
inside it was found to be correct and still unable to prove what a society would need to know.
Credential authenticity is not current authorization. Current authorization is not privacy. Privacy
is not freedom from coercion. Federation is not institutional independence. A transparency log is not
protection from a split view. And a green test is not proof that the test asked the right question.
Polaris is the accumulation of those sentences into mechanisms, and the honest map of where the
accumulation stands: at version 9.345 for the body of this paper, and at version 9.355 for the
holder-side sections, which were rewritten when the work they described as missing was built.
